\section{Preliminaries}
\label{sec:preli}
For database vector $x$ and query $q$ in $\mathbb{R}^d$, JHQ applies a square
orthogonal transform $Q$: $y=Qx$ and $q'=Qq$. It preserves Euclidean distance
without reducing dimension. The transformed vector is split into $M$ subspaces
of dimension $D_s$, with $d=MD_s$. Each $B$-bit primary code selects one of
$2^B$ codewords formed as Cartesian products of one-dimensional Lloyd--Max
levels. We use $B=D_s=8$: one bit per coordinate and one byte per subspace.

The residual is $r=y-\hat{y}_{\mathrm{primary}}$, measured against the primary
reconstruction, not an IVF centroid. With $B_r=8$ residual bits per dimension,
the logical representation stores one primary and eight residual bits per
dimension. IVF restricts the candidate set without changing this representation.

For every routed candidate, query-specific lookup tables evaluate the primary
score $D_P$; the query itself is not quantised into a database code. JHQ retains
$c_k$ primary survivors, evaluates their hierarchical score $D_H$, and returns
the top-$k$ under $D_H$. We write $c_k=\alpha k$; the implementation uses
$\max(k,\lceil\alpha k\rceil)$. Refinement cannot recover neighbours excluded
by routing or primary selection.
